\documentclass[11pt]{article}
\usepackage[a4paper,margin=1in]{geometry}
\usepackage[T1]{fontenc}
\usepackage[utf8]{inputenc}
\usepackage{lmodern}
\usepackage{amsmath,amssymb,mathtools}
\usepackage{microtype}
\setlength{\parindent}{1.2em}
\setlength{\parskip}{0.25em}
\newcommand{\dd}{\,d}
\newcommand{\modulo}[1]{\pmod{#1}}
\begin{document}

\begin{center}
{\Large G. Lejeune Dirichlet's Werke. II.}\par\vspace{1.2em}
{\Large XVIII.}\par\vspace{0.5em}
{\LARGE \textbf{Sur l'équation \(t^2+u^2+v^2+w^2=4m\)}}\par\vspace{1em}
{\large Par M. G. Lejeune Dirichlet}\par\vspace{0.5em}
{\small Liouville, Journal de Mathématiques pures et appliquées, Série II, Tome I, 1856, p. 210--214.}
\end{center}

\bigskip
\begin{center}
{\large \textbf{Sur l'équation \(t^2+u^2+v^2+w^2=4m\).}}
\end{center}

\begin{center}
[Extrait d'une Lettre de M. Lejeune Dirichlet à M. Liouville.]
\end{center}

Permettez-moi, cher ami, de revenir un instant sur la conversation que nous avons eue dernièrement sur le beau théorème de Jacobi relatif au nombre des décompositions d'un entier en quatre carrés, théorème que l'illustre géomètre a tiré d'abord de ses séries elliptiques et dont il a donné depuis une démonstration arithmétique\footnote{Journal de Crelle, tome XII, page 167; Jacobi's gesammelte Werke, Bd. VI, S. 245.}. Ayant mieux rappelé mes souvenirs, je vais vous présenter avec plus de développement, et surtout avec un peu plus d'ordre, ce que l'autre jour je n'ai pu que vous indiquer. Comme je vous le disais, j'ai toujours éprouvé quelque difficulté, ce qui, je crois, est arrivé aussi à d'autres personnes, lorsque j'ai voulu reproduire cette belle démonstration sans avoir le texte de l'auteur sous les yeux et sans recourir aux opérations sur les séries, dont la démonstration de Jacobi, ainsi qu'il en avertit lui-même, n'est que la traduction. Cette circonstance m'a fait chercher il y a déjà longtemps à fonder sur d'autres principes une nouvelle démonstration du théorème dont il s'agit, mais je n'en parlerai pas ici, cette démonstration devant trouver naturellement sa place dans un travail dont la rédaction m'occupe depuis quelque temps. Pour le moment, je n'ai d'autre objet que de simplifier celle de Jacobi, et de la rendre surtout plus facile à retenir, en mettant dans tout son jour le fait arithmétique ou plutôt algébrique qui en forme le principal fondement.

Pour éviter des répétitions inutiles, je remarquerai que tous les entiers que nous désignerons par des lettres latines sont positifs et impairs, à l'exception de ceux désignés par \(x\) et \(x'\), qui ne sont pas assujettis à cette limitation.

Quant au lemme qui sert de point de départ à Jacobi et qui définit le nombre des solutions de l'équation:
\[
t^2+u^2=2p,
\]
\(p\) étant supposé donné, je ne reproduirai pas ici la démonstration très simple que Jacobi en a donnée, d'autant plus que ce lemme rentre dans une proposition générale dont j'ai eu à parler ailleurs\footnote{Journal de Crelle, tome XXI, page 3; Bd. I, S. 468 dieser Ausgabe von G. Lejeune Dirichlet's Werken. K.}. On sait qu'il faut opérer toutes les décompositions:
\[
p=ad
\]
de \(p\) en deux facteurs, et que le nombre en question est égal à l'excès du nombre des cas où le premier facteur \(a\) est de la forme \(4\nu+1\), sur celui des cas où \(a\) est de la forme \(4\nu+3\); en convenant donc de faire correspondre à chaque décomposition une unité \(\delta\), positive ou négative selon ces deux cas, le nombre dont il s'agit sera exprimé par la somme:
\[
\sum \delta.
\]

Cela posé, cherchons à déterminer le nombre des solutions de l'équation:
\[
\tag{1} t^2+u^2+v^2+w^2=4m,
\]
dans laquelle \(m\) indique un entier impair donné. Pour obtenir toutes ces solutions, il faudra poser de toutes les manières possibles:
\[
4m=2p+2q\quad\text{ou}\quad 2m=p+q,
\]
et résoudre ensuite pour chaque couple \(p,q\), de la manière la plus générale, les deux équations:
\[
t^2+u^2=2p,\qquad v^2+w^2=2q.
\]
Si, considérant d'abord un couple déterminé \(p,q\), nous faisons correspondre à chaque décomposition:
\[
q=bc,
\]
une unité \(\varepsilon=\pm1\), qui dépende de \(b\) de la même manière que \(\delta\) de \(a\), le nombre des solutions de l'équation (1) qui répondent à ce couple \(p,q\), sera:
\[
\sum\delta\cdot\sum\varepsilon,
\]
ou encore:
\[
\sum\eta,
\]
les termes de cette somme répondant aux combinaisons différentes que l'on peut former avec les décompositions:
\[
p=ad,\qquad q=bc,
\]
et \(\eta\) étant égal à l'unité positive ou négative, selon que les restes de \(a\) et \(b\), lorsqu'on les divise par 4, sont ou ne sont pas égaux, ou, ce qui revient au même, selon que \(a-b\) est ou n'est pas divisible par 4. Ceci bien entendu, il est manifeste que le nombre de toutes les solutions de notre équation (1) sera encore exprimé par la somme:
\[
\sum\eta,
\]
dont les termes dépendront toujours de la même manière de \(a\) et \(b\), mais répondront maintenant à toutes les combinaisons que \(a\) et \(b\) présentent dans la solution complète de l'équation:
\[
\tag{2} 2m=ad+bc.
\]

Avant d'aller plus loin, il est bon de faire une remarque très simple et qui nous sera utile plus tard: c'est que, parmi les deux entiers pairs:
\[
a-b\quad\text{et}\quad c+d,
\]
il y en a toujours un, et qu'il n'y en a qu'un qui soit divisible par 4. En effet, s'il en était autrement, on aurait
\[
\text{ou } a\equiv b,\ c\equiv -d,\quad
\text{ou } a\equiv -b,\ c\equiv d \pmod{4},
\]
et il en résulterait \(ad+bc\equiv0\pmod{4}\). On peut donc définir \(\eta\) autrement en disant que \(\eta\) est égal à \(-1\) ou à \(+1\), selon que \(c+d\) est ou n'est pas divisible par 4.

Distribuons maintenant les solutions de l'équation (2) en deux classes, en comprenant dans la première celles qui sont telles que \(a=b\), et dans la seconde toutes les autres. Comme ces dernières sont associées deux à deux et résultent l'une de l'autre, en échangeant à la fois \(a\) avec \(b\), et \(d\) avec \(c\), ce qui ne changera pas la valeur de \(\eta\), on voit que le nombre \(H\) des solutions de l'équation (1) peut être mis sous la forme:
\[
H=\sum\eta' + 2\sum\eta'',
\]
les termes \(\eta'\) de la première somme, qui sont d'ailleurs tous égaux à l'unité positive, répondant aux solutions de l'équation (2) pour lesquelles on a \(a=b\), et ceux \(\eta''\) de la seconde aux solutions de la même équation qui satisfont à la condition \(a>b\).

Cela posé, nous allons considérer d'abord les solutions (2) pour lesquelles on a \(a>b\). Posant pour cela:
\[
a'=c(x+1)+d(x+2),\qquad c'=a(x+1)-b(x+2),
\]
\[
b'=cx+d(x+1),\qquad d'=-ax+b(x+1),
\]
nous aurons identiquement:
\[
a'd'+b'c'=ad+bc.
\]
Il résulte de là que toute solution de l'équation (2) en fournit une infinité d'autres au moyen de l'entier indéterminé \(x\), mais il reste à voir à quelle limitation cet entier doit être assujetti pour que \(a',b',c',d'\) soient impairs et positifs, et qu'on ait de plus \(a'>b'\) comme dans la solution primitive. Or chacun des couples:
\[
x,\ x+1\quad\text{et}\quad x+1,\ x+2,
\]
étant composé de deux entiers, l'un pair, l'autre impair, la première condition se trouve toujours remplie. Quant à la seconde, comme on a:
\[
c'=(a-b)(x+1)-b,\qquad d'=b-(a-b)x,
\]
on voit qu'elle détermine complètement l'entier \(x\), \((a-b)x\) devant être le multiple du nombre pair et positif \(a-b\), immédiatement inférieur à l'impair \(b\), et l'on voit encore, \(x\) n'étant pas négatif, que \(a'\) et \(b'\) seront positifs, et qu'on a de plus \(a'>b'\).

Après nous être assurés que les expressions données plus haut fournissent toujours une solution unique assujettie aux mêmes conditions que la solution primitive, voyons à quelle solution nous serons conduits, si nous prenons la solution \(a',b',c',d'\) à son tour pour point de départ. Comme pour cela il faut dans les expressions:
\[
c'(x'+1)+d'(x'+2),\qquad a'(x'+1)-b'(x'+2),
\]
\[
c'x'+d'(x'+1),\qquad -a'x'+b'(x'+1),
\]
donner à l'indéterminée \(x'\) la valeur entièrement déterminée en vertu de ce qui précède, qui rend ces expressions positives et que, d'un autre côté, les équations posées plus haut donnent par leur résolution celles-ci:
\[
a=c'(x+1)+d'(x+2),\qquad c=a'(x+1)-b'(x+2),
\]
\[
b=c'x+d'(x+1),\qquad d=-a'x+b'(x+1),
\]
dont les premiers membres sont positifs, on voit que \(x'\) coïncide avec \(x\), et que nous nous trouvons ramenés à la solution primitive \(a,b,c,d\).

Si maintenant nous observons qu'à nos deux solutions ainsi liées entre elles, et qui se déduisent mutuellement l'une de l'autre par un procédé uniforme, correspondent toujours des valeurs opposées de \(\eta''\), comme cela résulte de l'équation:
\[
a-b=c'+d'
\]
et de la remarque faite plus haut, nous concluons que \(\sum\eta''\) se réduit à zéro, et nous aurons simplement:
\[
H=\sum\eta'.
\]
Pour évaluer cette somme dont tous les termes \(\eta'\) sont égaux à l'unité positive, tout se réduit à trouver le nombre des solutions de l'équation:
\[
a(d+c)=2m.
\]
Or \(a\) étant un diviseur déterminé de \(m\), \(m/a=e\) sera aussi un tel diviseur, et l'on aura à résoudre l'équation:
\[
d+c=2e,
\]
dont le nombre des solutions est \(e\); d'où il suit, \(a\) et par conséquent aussi \(e\) devant être successivement égalés à tous les diviseurs de \(m\), que le nombre \(H\) qu'il s'agissait de déterminer coïncide avec la somme des diviseurs de \(m\).

\end{document}
